\chapter{Sum and Density Preserving Permutations}

\section{Agnew's Theorem}

Let's start this section proving a really famous theorem due to Agnew, which characterizes all sum preserving permutations

\begin{definition}{}{}
    A permutation $p\in\sym$ is called Agnew if there exists a $k\in\omega$ such that for all $n\in\omega$, $p[n]$ can be written as a union of at most $k$ intervals.
\end{definition}

Equivalently, we can say that $p$ is Agnew if we can switch between elements in $\text{Im}(p)$ and outside at most $k$ times. To formalize such condition, say $M<N$ when $\max(M)<\min(N)$, and call $M=\{m_0<m_1<\dots<m_k\}$ and $N=\{n_0<n_1<\dots<n_k\}$ collated sets if their elements alternate between each other, i.e., $m_0<n_0<m_1<\dots<m_k<n_k$.

\begin{definition}{}{}
    We say $p\in\sym$ satisfies condition A if there exists a $k\in\omega$ such that whenever $M$ and $N$ are collated sets and $p[M]<p[N]$, we have $|M|=|N|<k$.
\end{definition}

It's easy to see that $p$ is Agnew iff $p^{-1}$ satisfies condition $A$. We can now state one of the main Theorems of this section.

\begin{theorem}{}{}
    $p\in\sym$ is Agnew iff it's a sum preserving permutation.
\end{theorem}

\begin{sketch}
    $(\Rightarrow)$ Given any $s=\sum_na_n$, we know by Cauchy's Criterion that, varying $m$, the partial sums starting from $a_m$ can be made arbitrarily small. So we can choose $n$ such that $p[n]$ covers $[0,m]$, meaning all at most $k$ gaps between the intervals in the image of $p$ can be made arbitrarily small.

    $(\Leftarrow)$ Assume $p$ is not Agnew, we will find a $\sum_na_n$ whose rearrangement corrupts. Since $p$ is not Agnew, we can find arbitrarily large alternating sequences of elements in $\text{Im}(p)$ and outside of it.

    Fix two such elements relative to $p[n_1]$ and associate the one inside $\text{Im}(p)$ with $1$, and the other with $-1$. Now, let $n_2$ such that $p[n_2]$ covers the previous interval and guarantees us to have four new alternating elements. Associate the ones inside $\text{Im}(p)$ with $\frac12$, and the ones outside with $-\frac12$. In general, choose $n_k$ such that $p[n_k]$ covers all previous elements and guarantees us $2k$ new elements, half inside and half outside of $\text{Im}(p)$. Associate the ones outside with $\frac1k$, and the ones inside with $-\frac1k$. Now, for all the remaining indices, just let them to be $0$, this defines a sequence $a_n$, where $a_n$ is the value we associate to its index $n$. Clearly for any $k$, we can let $\sum a_n$ be such that all its partial sums eventually has value at most $\frac1k$, meaning it converges to $0$. However, by construction, the partial sums of the rearranged series up to $n_k$ is always $1$.
\end{sketch}

Let's now write all these steps more formally

\begin{proof}
    $(\Rightarrow)$ By Cauchy Criterion, given $\varepsilon>0$, there is an $m_0\in\omega$ such that $\forall n>m\geq m_0$, we have
    $$\left\lvert\sum_{i=m}^na_i\right\rvert<\frac{\varepsilon}{2k}$$
    Let then $n$ be large enough such that $p[n]$ covers $[0,m_0]$, then, since $p$ is Agnew, we can write $p[n]=I_1\cup I_2\cup\dots\cup I_\ell$, where $\ell\leq k$ and $I_i=[b_i,B_i]$. Therefore, since any $m\in I_j$ has $m\geq m_0$, for $j=1,\dots,\ell$, then, by Cauchy's Criterion
    $$\left\lvert\sum_{i=0}^na_{p(i)}-\sum_{i=0}^{b_1}a_i\right\rvert\leq\left\lvert\sum_{i\in I_1}a_i\right\rvert+\dots+\left\lvert\sum_{i\in I_\ell}a_i\right\rvert<\frac\varepsilon2$$
    Therefore, we have that
    $$\left\lvert\sum_{i=0}^na_{p(i)}-\sum_{i=0}^{n}a_i\right\rvert\leq\left\lvert\sum_{i=0}^na_{p(i)}-\sum_{i=0}^{b_1}a_i\right\rvert+\left\lvert\sum_{i=0}^{b_1}a_i-\sum_{i=0}^na_i\right\rvert<\varepsilon$$

    $(\Leftarrow)$ \textcolor{laser_eyes}{É bem mais fácil de ver com um desenho colorido :)}
\end{proof}

It would be really nice if condition A also characterizes density preserving permutations. Unfortunately, this is not the case, we will see later that condition A preserves what we call strong density, or uniform density.

\section{Relative Density Preserving Permutations}

Now, as expected, the previous condition would not be called condition A if we didn't have a condition B:

\begin{definition}{}{}
    We say $p\in\sym$ satisfies condition B if there exists a $k\in\omega$ such that whenever $|M|=|N|$, with $M<N$ and $p[N]<p[M]$, we have $|M|=|N|<k$.
\end{definition}

It's not hard to see that condition B implies condition A, but the converse is not true. The following lemma actually shows something stronger:

\begin{lemma}{}{}
    Condition B implies condition A, but we can find some $p$ such that both $p$ and $p^{-1}$ satisfies A, but not B.
\end{lemma}

\begin{proof}
    Consider $M$, $N$ collated sets, i.e., $m_1<n_1<m_2<n_2<\dots<m_j<n_j$, such that $p[M]<p[N]$, then, if $\ell$ is like a half of $j$ (you could take the floor), then
    $$n_1<n_2<\dots<n_\ell<m_{\ell+1}<\dots<m_j$$
    or $j-1$ so both sets $N'=\{n_1,\dots,n_\ell\}$ and $M'=\{m_{\ell+1},\dots,m_j\}$ have the same size. Therefore, $N'<M'$, but $p[M']<p[N']$, then, by condition $B$, $\lvert N'\rvert=\ell<k$, meaning $j$ is at most $2k$ (or $2k+1$ depending on the parity).

    To construct the counterexample, let $J_1=\{1\}$, $J_2=\{2,3\}$, $J_3=\{4,5,6\}$, etc. And consider the unique $p$ that biject $J_k$ in $J_k$ reserving the order of its elements. Then, for any $k$, we can take $M$ as the first half of $J_{2k}$ and $N$ as the remainder, then $M<N$, but $p[N]<p[M]$. However, note that $p=p^{-1}$ and $p$ is clearly Agnew, since the image is the union of at most 2 intervals.
\end{proof}

It would then be expected, based on the discussions until now, that condition B would characterize the permutations that preserve relative density. However, this is also not the case:

\begin{theorem}{}{}
    $p\in\sym$ satisfies condition B iff it preserves relative density.
\end{theorem}

\begin{proof}
    $(\Rightarrow)$ For any $n$, let $M=\{m<n:p(m)\geq n\}$ then, since $p$ is a permutation, we can find some $N\subseteq\{m\geq n:p(m)<n\}$ with the same size as $M$. Then $M<N$, but $p[N]<p[M]$, so $\lvert M\rvert=\lvert N\rvert<k$, which implies that, if $A\subseteq\omega$ is infinte coinfinite with density, then
    $$\lvert\lvert p[A]\cap n\rvert-\lvert A\cap n\rvert\rvert<k\implies \lvert p[A]\cap n\rvert=\lvert A\cap n\rvert+m_n^A$$
    with $\lvert m_n^A\rvert<k$. This shows $p$ preserves density. Now, given $A\subseteq B\subseteq \omega$ infinite and coinfinite with relative density, by the previous result we have that
    $$\frac{\lvert p[A]\cap n\rvert}{\lvert p[B]\cap n\rvert}=\frac{\lvert A\cap n\rvert+m_n^A}{\lvert B\cap n\rvert+m_n^B}=\frac{\frac{\lvert A\cap n\rvert}{\lvert B\cap n\rvert}+\frac{m_n^A}{\lvert B\cap n\rvert}}{1+\frac{m_n^B}{\lvert B\cap n\rvert}}\to d_B(A)$$
    then $d_{p[B]}(p[A])=d_B(A)$, i.e., it preserves relative density.

    $(\Leftarrow)$ \textcolor{laser_eyes}{This one is trickier.}
\end{proof}

What does such result reveal to us? Not that much, to be fair. Actually, it just says that the set of permutations that preserve relative density are also sum-preserving, but it says nothing about the set of permutations that preserve density, which is what we want.

\section{Lévy's Group}

Before we proceed, look again at the $(\Rightarrow)$ case in the previous proof, we showed that condition B implies the following condition:
\begin{equation}\tag{L}
    \lim_{n\to\infty}\frac{\lvert\{m\in\omega:m<n\leq p(m)\}\rvert}{n}=0
\end{equation}
such condition L (due to Lévy) has some interesting properties that was deeply studied.

\begin{definition}{}{}
    The set $\mc{G}$ of permutations that satisfies L form a group, known as Lévy's Group.
\end{definition}

\begin{proof}
    \textcolor{laser_eyes}{Pendente}
\end{proof}

And, in relation to what it preserves, we have the following theorem:

\begin{theorem}{}{}
    $p\in\mc{G}$ if, and only if, for all $A\subseteq\omega$
    $$\lim_{n\to\infty} d_n(A)-d_n(p[A])=0$$
\end{theorem}

\begin{proof}
    $(\Rightarrow)$ is a scholium of the start of the previous proof;

    $(\Leftarrow)$ if $p\notin\mc{G}$, choose $(a_n)$ fast growing such that $\frac{a_{n-1}}{a_n}\to0$ and $E_n=\{m\in\omega:m<a_n\leq p(m)\}$ satisfies $\frac{\lvert E_n\rvert}{a_n}>\alpha$, for some $\alpha>0$. Then, let $A=\bigcup_nE_n$. We have that
    $$d_{a_n}(A)\geq d_{a_n}(E_n)>\alpha>0$$
    and, since $p[E_j]\geq n_j$, we have $p[A]\cap a_n\subseteq\bigcup_{j<n}p[E_j]$, so
    $$d_{a_n}(p[A])\leq d_{a_n}\rp{p\left[\bigcup_{j<n}E_j\right]}=d_{a_n}\rp{\bigcup_{j<n}E_j}\leq\frac{a_{n-1}}{a_n}\to0$$
    then $\lim_{n\to\infty}d_{a_n}(A)-d_{a_n}(p[A])\geq\alpha>0$.
\end{proof}

i.e., the permutations that satisfies L ``preserves'' density in all cases, including when $d(A)=\text{osc}$.

This shows that, if we had a condition C such that $p$ satisfies C iff $p$ preserves density, then condition L would imply it. Let's show such implications are strict.

\begin{theorem}{}{}
    $\text{B}>\text{L}>\text{C}$ and $\text{A}\perp\text{L},\text{C}$.
\end{theorem}

\begin{proof} We already know that $\text{B}\geq\text{L}\geq\text{C}$, then it just remains to show the strict inequalities and the incomparability with condition A.
    \begin{itemize}
        \item $\text{A}\not\geq\text{L}:$ Let $N_{k+1}\gg N_k$, and consider $p$ that swaps $[N_k,2N_k)$ with $[2N_k,3N_k)$, for each $k\in\omega$. Clearly $p$ satisfies condition $A$, since $p^{-1}=p$ is Agnew (the image $p[n]$ is always the union of at most two intervals). However, using the notation of the previous proof, we have $\lvert E_{2N_k}\rvert=N_k$, then it doesn't satisfy condition L.
        
        \item $\text{C}\not\geq\text{L}$: The permutation defined in the previous item works, since
        \begin{align*}
            \lvert p[A]\cap n\rvert & = \lvert\{k<n:p(k)<n\}\rvert+\lvert\{k\geq n:p(k)<n\}\rvert\\
            & = \lvert A\cap n\rvert-\lvert\{k<n:p(k)\geq n\}\rvert+\lvert\{k\geq n:p(k)<n\}\rvert
        \end{align*}
        

        \item $\text{L}\not\geq\text{B}:$ We now need to swap two chunks, but whose size grows sub linearly in some sense. Let then $M_k=[2^k,2^k+k)$ and $N_k=[2^k+k,2^k+2k)$ and $p$ be the unique permutation that swap both, for each $k\in\omega$. Then, clearly $p$ doesn't satisfy condition $B$, however, for $2^k\leq n<2^{k+1}$, the maximum amount $E_n$ can attain is at $n=2^k+k$, where
        $$\frac{\lvert E_n\rvert}n=\frac{k}{2^k+k}\to0$$
        so $p$ does satisfy condition L.

        \textcolor{laser_eyes}{It's worth noting that this proof was unnecessary: if $\text{L}\geq\text{B}$, since $\text{B}\geq\text{A}$, we would have $\text{L}\geq\text{A}$, contradicting the next item.}
    
        \item $\text{L}\not\geq\text{A}:$ The heuristic is clear, based on the previous one. We just let $J_k=[2^k,2^k+2k)$, and let $p$ be such that the even numbers in $J_k$ are mapped to its first $k$ elements, and the odd ones to the last $k$. Then, by the same argument as before, $p$ satisfies condition L, but we have arbitrarily large collated sets that get separated, so it doesn't satify condition A.

        \item $\text{C}\not\geq\text{A}$: Since $\text{L}\geq\text{C}$, if $\text{C}\geq\text{A}$, then $\text{L}\geq\text{A}$, which is false by the previous item.

        \item $\text{A}\not\geq\text{C}$: Let $a_0=0$ and $a_{k+1}=2^{a_k}$, let $p$ be such that it swaps the interval $I_k=[a_k,2a_k)$ with $J_k=[a_k^2,a_k^2+a_k)$, and let $A=2\mbb{N}\setminus\bigcup_k J_k$. Then, when $n\in J_k$, we have that
        $A\cap n$ has all the even elements until $n$, but at most the $a_k$ ones in $I_k$. But, since $n\geq a_k^2$, at this point these missed elements are negligible, and then $d(A)=\frac12$.

        However, for $n=2a_k$, we have $p[A]\cap n$ missing all $a_k$ elements in $I_k$, then, since $n=2a_k$ and $p[A]$ are the even naturals with gaps, it has at most $\frac{a_k}{2}$ elements in it, i.e., $d_n(p[A])\leq\frac14$.
    \end{itemize}
\end{proof}
